\documentclass[article,12pt]{elsarticle}
\usepackage[spanish]{babel}
\usepackage[utf8]{inputenc}
\usepackage{amssymb}
\usepackage{flushend}
\usepackage{float}
\usepackage{anysize}
\usepackage{amsmath}
\usepackage{amsfonts}
\usepackage{dcolumn}
\usepackage{amsthm}
\usepackage{caption}
\usepackage{float}
\usepackage{subfigure}
\usepackage{graphicx}
\usepackage{lmodern}
\usepackage{booktabs}
\usepackage{color}
\usepackage{longtable}
\usepackage{rotating}
\usepackage{etoolbox}
\usepackage{comment}
\usepackage{pdflscape}
\usepackage{booktabs}


\begin{document}

\begin{frontmatter}
    \title{Laboratorio de Fenómenos Colectivos \\  Tarea de tensión Superficial} \\
    \author{Andrés López González$^1$} %Aquí va el nombre de los integrantes que participaron en el informe.
    \address{$^1$ Facultad de Ciencias, Universidad Nacional Autónoma de México, M\'exico CDMX.}
\end{frontmatter}

La diferencia en el comportamiento térmico del agua en un recipiente abierto frente a uno cerrado, así como el impacto del material del recipiente, puede explicarse mediante conceptos de **transferencia de calor**, **evaporación** y **propiedades térmicas de los materiales**. 

### 1. **Caso de un recipiente abierto: por qué el agua está más fría que la temperatura ambiental**
Cuando el recipiente está abierto, ocurren los siguientes fenómenos:

#### a) **Evaporación**
El agua constantemente pierde energía térmica debido al proceso de evaporación. Este proceso ocurre porque las moléculas en la superficie del agua con mayor energía cinética escapan al aire como vapor. Esto genera un enfriamiento del líquido restante porque se pierden las moléculas con mayor energía, reduciendo la temperatura promedio del agua.

La tasa de evaporación depende de:
- **La temperatura del agua y del aire.** Si el aire está seco (baja humedad relativa), la evaporación será más rápida, y el agua se mantendrá más fría.
- **La velocidad del viento.** Una corriente de aire aumenta la evaporación al remover las moléculas de vapor cercanas a la superficie del agua.
- **La exposición al aire libre.** Cuanto mayor sea la superficie del agua expuesta, mayor será la tasa de evaporación.

#### b) **Balance energético con el ambiente**
En un recipiente abierto, el intercambio de calor con el aire ambiente no se limita a la conducción o convección, sino que está influenciado también por la pérdida de calor por evaporación. Aunque el agua recibe calor del aire y del recipiente (por conducción y convección), la evaporación puede superar esta ganancia energética, manteniendo la temperatura del agua más baja que la del aire.

---

### 2. **Caso de un recipiente cerrado: por qué el agua alcanza la temperatura ambiental**
Cuando el recipiente está cerrado (pero no adiabático), se elimina la posibilidad de evaporación. En este caso, los únicos mecanismos de transferencia de calor son:
- **Conducción:** Calor que fluye desde las paredes del recipiente al agua.
- **Convección:** Transferencia de calor por el movimiento del aire dentro del recipiente.
  
En ausencia de evaporación, el sistema alcanza un **equilibrio térmico**, donde el agua y el aire dentro del recipiente logran la misma temperatura. Este proceso es más rápido si las paredes del recipiente permiten una buena transferencia de calor.

---

### 3. **Influencia del material del recipiente**
El material del recipiente afecta la transferencia de calor y, por ende, la temperatura del agua. Veamos las propiedades térmicas de los materiales mencionados:

#### a) **Metálico**
- **Alta conductividad térmica:** El metal transfiere calor de manera muy eficiente. Si el recipiente está abierto, esta eficiencia puede acelerar el calentamiento del agua desde el ambiente, pero el efecto de enfriamiento por evaporación seguirá presente. En un recipiente cerrado, el agua alcanzará rápidamente la temperatura ambiental.

#### b) **Vidrio**
- **Conductividad térmica moderada:** El vidrio transfiere calor menos eficientemente que el metal, por lo que los cambios de temperatura son más lentos. Esto puede hacer que el agua en un recipiente abierto pierda calor ligeramente más lento, pero la evaporación seguirá predominando.

#### c) **Unicel (poliestireno expandido)**
- **Baja conductividad térmica:** El Unicel es un excelente aislante térmico. En un recipiente abierto, la evaporación será el principal mecanismo de pérdida de calor, ya que la transferencia de calor entre el agua y el ambiente es mínima. En un recipiente cerrado, el agua tardará mucho más en alcanzar la temperatura ambiental debido a la baja eficiencia de intercambio térmico.

---

### 4. **Diferencias en la temperatura del agua según el material**
En un recipiente abierto:
- En **metal**, el agua puede estar ligeramente más cálida que en vidrio o Unicel debido a la mejor transferencia de calor del metal hacia el agua.
- En **vidrio**, la temperatura del agua será intermedia.
- En **Unicel**, el agua puede estar más fría, ya que el aislamiento térmico limita la ganancia de calor desde el ambiente, mientras que la evaporación sigue enfriando el agua.

En un recipiente cerrado:
- No habrá diferencias significativas en la temperatura final alcanzada, pero los tiempos para llegar al equilibrio térmico serán más cortos en recipientes metálicos y más largos en recipientes de Unicel.

\end{document}

